% !TEX root = ../notes.tex

\documentclass[../notes.tex]{subfiles}

\begin{document}

\section{March 11}
We now turn to archimedean theory.

\subsection{Real Lie Groups}
We would like to retell our nonarchimedean story for real Lie groups. Thus, we want to define smooth representations, produce a Hecke algebra, and define a notion like spherical representations.

Let's start by explaining what real Lie groups we will consider, and we will discuss some of their structure theory.
\begin{definition}[Lie group]
	A \textit{Lie group} is a group object in the category of smooth manifolds.
\end{definition}
\begin{example}[classical groups]
	The classical groups look like $\op{GL}_n(\RR)$, $\op O(p,q)$, $\op{GL}_n(\CC)$, $\op U(p,q)$, $\op{GL}_n(\HH)$, $\op{Sp}(p,q)$, $\op{Sp}_{2n}(\RR)$, and a few others. All of these are real points of real algebraic groups. For example, $\HH$ embeds into $M_4(\RR)$, so $\op{GL}_n(\HH)$ embeds into $\op{GL}_{4n}(\RR)$.
\end{example}
The plethora of examples is more or less because we have chosen to study our groups over $\RR$ instead of over $\CC$.
\begin{remark}
	For any real algebraic group $G$, one can base-change to $\CC$, and then $G_\CC$ is classified by the root datum (by \Cref{thm:root-datum}).
\end{remark}
\begin{example}
	All the groups $\op{GL}_n(\RR)$, $\op U(p,q)$ with $p+q=n$, and $\op{GL}_{n/2}(\HH)$ all base-change to $\op{GL}_{n,\CC}$.
\end{example}
Lie groups come with a distinguished compact subgroup.
\begin{theorem}[Cartan]
	Fix a connected real Lie group $G$. Then $G(\RR)$ admits a maximal compact subgroup, and all maximal compact subgroups are conjugate.
\end{theorem}
\begin{example}
	The maximal compact subgroup $\op{GL}_n(\RR)$ is $\op O(n)$. The maximal compact subgroup of $\op U(p,q)$ is $\op U(p)\times\op U(q)$.
\end{example}
\begin{example}
	It turns out that $\op{Sp}_{2n}(\RR)$ admits $\op U(n)$ as a maximal compact subgroup. To explain the embedding, note that one can give $\RR^{2n}$ a complex structure so that the standard symplectic form $\omega$ is the imaginary part of a positive-definite Hermitian form $h$ on $\CC^n$. Then $\op U(h)$ embeds into $\op{Sp}(\omega)$.
\end{example}
\begin{remark}
	These maximal compact subgroups are important because they (roughly) play the role of our $G(\OO)$ in our nonarchimedean theory. For example, the name ``spherical vector'' actually comes from this archimedean theory: a spherical vector for $\op{GL}_n(\RR)$ should be fixed by the rotation subgroup $\op O_n(\RR)$!
\end{remark}
Let's explain how we found so many real groups.
\begin{theorem}[Cartan]
	Fix a complex reductive group $G$ over $\CC$.
	\begin{listalph}
		\item There are finitely many real algebraic groups $G_0$ for which $(G_0)_\CC$ (up to isomorphism).
		\item The set of isomorphism classes is in bijection with Cartan involutions $\theta\colon G\to G$ (satisfying $\theta^2=1$), up to conjugacy by automorphisms.
	\end{listalph}
\end{theorem}
Roughly speaking, the involution $\theta$ should be treated as complex conjugation.
\begin{example}
	Consider the group $G=\mathbb G_{m,\CC}$. There are two involutions $\theta$, which are the identity and the inversion. The identity corresponds to the real form $\op U(1)$, and the inverse corresponds to $\mathbb G_{m,\RR}$.
\end{example}
\begin{example}
	Consider the group $T=\mathbb G_{m,\CC}^r$. Then $\theta$ should be chosen among automorphisms of $X_*(T)$. If one can diagonalize the action of $\theta$ on the lattice $X_*(T)$, then it diagonalizes into some copies of $(+1)\id$ and $(-1)\id$, so the real form is a product of $\op U(1)$s and $\mathbb G_{m,\RR}$s. However, there are more real forms: for example, $\op{Res}_{\CC/\RR}\mathbb G_{m,\CC}$ is a real form of $\mathbb G_{m,\CC}^2$, and it corresponds to the involution $\begin{bsmallmatrix}
		0 & 1 \\ 1 & 0
	\end{bsmallmatrix}$ on the lattice.
\end{example}
\begin{example}
	Consider the group $G=\op{GL}_{n,\CC}$. There are inner involutions, which look like conjugation by
	\[\op{diag}(\underbrace{+1,\ldots,+1}_p,\underbrace{-1,\ldots,-1}_q).\]
	These correspond to the real forms $\op U(p,q)$. However, $\op{GL}_{n,\CC}$ also has an outer automorphism given by $g\mapsto g^{-\intercal}$; equivalently, we may view this as a composite $\op{GL}_n(\CC)\cong\op{GL}(V)\cong\op{GL}(V^*)\cong\op{GL}_n(\CC)$, where the middle isomorphism is given by a perfect pairing $V\times V\to\CC$, where $\dim V=n$. (Alternatively, up to inner automorphisms, this outer automorphism is basically given by ``flipping'' the Dynkin diagram, which is the only nontrivial graph isomorphism.) There are basically two cases for such a perfect pairing.
	\begin{itemize}
		\item Symmetric: this always exists, and it gives the ``split'' form $\op{GL}_{n,\RR}$.
		\item Alternating: this exists when $n$ is even, and it gives the real form $\op{GL}_{n/2}(\HH)$.
	\end{itemize}
\end{example}

\subsection{Smooth Representations}
Let $G$ be a real algebraic group. Our next task is to figure out what our representations are.
\begin{definition}
	Fix a real algebraic group $G$, and let $V$ be a topological complex vector space. Then a \textit{continuous representation} of $G$ is one for which the action map
	\[G\times V\to V\]
	is continuous. If $V$ is further a Hilbert space, and the action preserves the Hermitian inner product, then the representation is \textit{unitary}.
\end{definition}
\begin{remark}
	It is typical to require $V$ to at least be a Fr\'echet space.
\end{remark}
\begin{example}
	Suppose that $G(\RR)$ acts smoothly on some real manifold $X$, and let $L^2\left(X,\left|\omega_X\right|^{1/2}\right)$ be the collection of half-densities, which we recall is a real line bundle on $X$. Then $G$ acts continuously on $L^2\left(X,\left|\omega_X\right|^{1/2}\right)$, and the $G$-action preserves the canonical inner product
	\[\langle f,g\rangle\coloneqq\int_Xfg.\]
	Here, $fg$ is a global section of $\left|\omega_X\right|$ and can therefore be integrated over $X$. However, we remark that $G$ also acts continuously on $L^p(X)$ for any $p\ge1$, so there are many representations here!
\end{example}
To understand our representations, we must walk the roads lain by Harish--Chandra.
\begin{definition}[smooth]
	Fix a continuous complex representation $V$ of a real Lie group $G$. Then a representation of a real algebraic group $G$ on $V$ is \textit{smooth} if and only if, for any $v\in V$ and continuous $\ell\in V^*$, the matrix coefficient $g\mapsto\ell(gv)$ is a smooth function on $G$.
\end{definition}
\begin{remark}
	It is not hard to check that $V^{\mathrm{sm}}$ is a subspace of $V$.
\end{remark}
\begin{remark}
	It turns out that $V^{\mathrm{sm}}$ is dense in $V$ as long as $V$ is Fr\'echet. Indeed, for any $\mu\in C_c^\infty(G)\,dg$, one can check that the integral
	\[\mu*v\coloneqq\int_Ggv\,d\mu(g)\]
	makes sense in the Fr\'echet space $V$. It is not hard to check that the smoothness of $\mu$ implies smoothness of $\mu*v$. By having $\mu$ concentrate around the identity, one can show that $\mu*v\to v$, so we see that the smooth vectors are dense.
\end{remark}
\begin{remark}
	The Lie algebra $\mf g$ of $G$ acts on $V^{\mathrm{sm}}$: namely, for $X\in\mf g$, being Fr\'echet implies that the limit
	\[Xv\coloneqq\lim_{t\to0}\frac{\exp(tX)v-v}t\]
	exists in $V$; in fact, this is a smooth vector!
\end{remark}
\begin{defihelper}[$K$-finite] \nirindex{K-finite@$K$-finite}
	Fix a continuous complex representation $V$ of a real Lie group $G$, and choose a compact subgroup $K\subseteq G$. A vector $v\in V$ is \textit{$K$-finite} if and only if
	\[\dim_\CC\op{span}\{kv:k\in K\}<\infty.\]
\end{defihelper}
\begin{remark}
	A similar argument is able to show that $(V^{\mathrm{sm}})_{\mathrm{fin}}\subseteq V$ is smooth when $V$ is a Fr\'echet space.
\end{remark}
The moral is that $V^{\mathrm{sm}}_{\mathrm{fin}}$ admits an action by $\mf g$ and an action by $K$. It turns out that this is the correct amount of data to remember from a Fr\'echet representation of a real Lie group.

\end{document}